\chapter{Experimental Chronology}
\label{app:chronology}
This appendix records the development sequence rather than presenting
only the final model. Phase names follow the repository. Infrastructure
checks and interrupted runs are identified separately from completed
retrieval evaluations.

\begin{center}
\small
\captionsetup{hypcap=false}
\captionof{table}{Chronology of the preserved experimental programme.}
\label{tab:chronology}
\begin{tabular}{@{}>{\raggedright\arraybackslash}p{2cm}>{\raggedright\arraybackslash}p{5.3cm}>{\raggedright\arraybackslash}p{5.8cm}@{}}
\toprule
Stage & Purpose & Outcome or status\\
\midrule
Initial prototype & Weaviate ingestion, JSON exports, sample retrieval & First working data path; no comparative vector-store benchmark\\
Foundation & Evidence filtering, parallel paper processing, Qdrant migration & Five chunk levels, resumable stages, stable identities\\
Early smoke & Verify retrieval and output writing & Tiny and empty artifacts are execution diagnostics, not full validation results\\
Integrity & Audit counts, offsets, checkpoints, and missing records & Dated problems documented; later frozen artifacts define final populations\\
Metric and storage & Per-chunk diagnostics, schema v2, configuration-aware identities & Persistence and idempotency checks; joined F1 remains primary\\
Fixed baseline & Score all five sizes, top five within each paper & Fixed 40 strongest; offline oracle records per-question opportunity\\
Embedding router & Logistic regression and 64-unit MLP & MLP gain below adoption threshold; logistic router retained\\
Mixed sizes & Pool candidates, with and without overlap removal & Neither raw nor deduplicated mixture beats fixed 40\\
Phase 1 & Frozen Qwen, question-only numeric decision & Valid outputs, strong preference for size 10\\
\bottomrule
\end{tabular}
\end{center}
\newpage
\begin{center}
\small
\textbf{Table~\ref{tab:chronology} (continued)}\par\medskip
\begin{tabular}{@{}>{\raggedright\arraybackslash}p{2cm}>{\raggedright\arraybackslash}p{5.3cm}>{\raggedright\arraybackslash}p{5.8cm}@{}}
\toprule
Stage & Purpose & Outcome or status\\
\midrule
Phase 2 preflights & Smoke, tiny-overfit and training-loop checks & Tail-accumulation issue corrected; interrupted v1 not a final model\\
Phase 2 v2 & Numeric supervised fine-tuning & Better label accuracy, weaker retrieval, large-size collapse\\
Phase 2B-A & Uniform restricted-alias classification & More medium-size predictions; strongest Qwen-only retrieval result\\
Phase 2B-B & Effective-number class weighting & No small-class recovery; weaker macro F1 and retrieval than A\\
Phase 2C & Base model with five-class head & Higher macro F1 than numeric SFT, but lower accuracy\\
Phase 2D & Exact token-count prompt & Better classification than 2C, slightly lower retrieval F1\\
Phase 2E & Three learning rates, five epochs & Fifteen candidates; winner 5e-6 at epoch four\\
Phase 3A & Level and hierarchy features with linear models & Small hierarchy difference; heuristics retained as diagnostics\\
Phase 3B & Weighted XGBoost on both representations & Training-fold macro F1 selects full tree; mean retrieval improves, median falls\\
Phase 3C & Qwen logits or hidden states fused with tree features & Initial result limited by in-sample base features\\
Phase 3C-OOF & Reconstruct base features using paper-grouped held-out predictions & Less optimistic fusion estimates; fixed full-data refit\\
Phase 4 & Five retrieval-regret regressors & Best learned global retrieval; still below fixed 40\\
Phase 5A & Gold-overlap supervision for local tree decisions & Close to, but not better than, matched fixed-40 tree retrieval\\
Phase 5A probe & Return all descendants at predicted level & Higher recall, larger budget, much lower F1\\
Phase 5B v1 & Frozen Qwen context category from free generation & Invalid-output gate fails; no downstream result\\
Phase 5B-v2 & Training-only prompt repair and constrained token scoring & Valid but nearly constant feature; no demonstrated retrieval gain\\
\bottomrule
\end{tabular}
\end{center}

Preflight runs check gradients, tiny-set fitting, accumulation, and
recovery; they are not independent estimates of generalisation.
Source transfers and hash audits likewise preserve an experiment
rather than create a new model. Phase 3A resumed gRPC feature
extraction through REST after 1,844 training questions. Its durable
processing-time sum must not be mistaken for uninterrupted wall time.

\chapter{Configurations and Output Contracts}
\label{app:protocols}

\section{Common retrieval settings}
The common settings are the candidate sizes 10, 20, 40, 80, and 160;
GPT-2 tokenisation; fixed 1,536-dimensional
\texttt{text-embedding-3-small} vectors; cosine similarity; source-paper
filtering; and top-five retrieval for the primary protocols.
The metric configuration hash preserved in the main frozen comparisons is
\begin{quote}\small
\path{9a3022fd1c808f72ccbf3265fe6020593bb58bdd28aeb9025b8c4b735d669de8}.
\end{quote}
The metric name is \texttt{f1\_joined\_topk}; the metric version is
\texttt{qasper-token-prf-v2}. A change in normalisation, candidate
population, or top-$K$ is a new evaluation configuration.

The three target definitions are
\begin{itemize}
\item \path{oracle-f1-evidence-smaller-v1}: legacy retrieval-derived labels;
\item \path{oracle-evidence-length-gpt2-smaller-midpoint-v1}: total evidence length;
\item local gold-overlap labels constructed for question--tree pairs.
\end{itemize}
The last target is described by its construction in
Chapter~\ref{ch:data}; it must not be assumed identical to a global
label merely because both return one of the same five numbers.

\section{Model identities and training}
The post-trained Qwen model is pinned to revision
\begin{quote}\small
\path{2fc06364715b967f1860aea9cf38778875588b17},
\end{quote}
and the base model to
\begin{quote}\small
\path{dc7cdfe2ee4154fa7e30f5b51ca41bfa40174e68}.
\end{quote}
Revision pinning matters because a model name alone does not identify
an immutable set of weights or tokenizer files.

The common three-epoch Qwen training setup uses seed 42, bfloat16,
microbatch size four, eight accumulation steps, maximum sequence
length 128, learning rate $2\times10^{-5}$, weight decay 0.01,
cosine scheduling, 5\% warm-up, and clipping at 1. The full training
population yields 71 optimiser steps per epoch. The five-epoch grid
changes its learning rate and schedule horizon as documented in
Chapter~\ref{ch:method}.

The Phase 2E recorded environment includes Python 3.11.13,
PyTorch 2.8.0 with CUDA 12.8, and Transformers
\texttt{5.15.0.dev0} pinned to commit
\begin{quote}\small
\path{2ef79f87a02111f8b49a72fb7d0c86b5b0bf10b7}.
\end{quote}
These describe the preserved GPU run, not requirements to install
on a computer that only needs to compile the thesis.

% Generated by scripts/generate_results.py; do not hand-edit.
\begin{table}[htbp]
\centering\small
\caption{All 15 Phase 2E validation candidates.}
\label{tab:lrgrid}
\begin{tabular}{@{}lrrrrrr@{}}
\toprule
LR & Epoch & Step & CE loss & Accuracy & Macro F1 & Top-two\\
\midrule
5e-6 & 1 & 71 & 1.4161 & 0.3084 & 0.1952 & 0.6017\\
5e-6 & 2 & 142 & 1.4814 & 0.2619 & 0.1505 & 0.5390\\
5e-6 & 3 & 213 & 1.4226 & 0.2857 & 0.1841 & 0.5714\\
5e-6 & 4 & 284 & 1.3760 & 0.3485 & 0.2278 & 0.6190\\
5e-6 & 5 & 355 & 1.3720 & 0.3485 & 0.2256 & 0.6310\\
1e-5 & 1 & 71 & 1.3605 & 0.3929 & 0.1789 & 0.6299\\
1e-5 & 2 & 142 & 1.3958 & 0.2998 & 0.1634 & 0.6515\\
1e-5 & 3 & 213 & 1.3921 & 0.3387 & 0.2122 & 0.6342\\
1e-5 & 4 & 284 & 1.4580 & 0.3236 & 0.2154 & 0.6039\\
1e-5 & 5 & 355 & 1.4509 & 0.3279 & 0.2126 & 0.6071\\
2e-5 & 1 & 71 & 1.3393 & 0.4091 & 0.1672 & 0.6537\\
2e-5 & 2 & 142 & 1.3556 & 0.3258 & 0.1729 & 0.7024\\
2e-5 & 3 & 213 & 1.5105 & 0.3463 & 0.2168 & 0.5909\\
2e-5 & 4 & 284 & 2.1367 & 0.3128 & 0.2232 & 0.5584\\
2e-5 & 5 & 355 & 2.2110 & 0.3171 & 0.2252 & 0.5563\\
\bottomrule
\end{tabular}
\par\smallskip\begin{minipage}{\textwidth}\footnotesize The selected candidate is 5e-6, epoch 4, step 284. The table does not imply retrieval was evaluated for every checkpoint.\end{minipage}
\end{table}


\section{Intermediate candidates and secondary variants}
The following tables retain the smaller comparisons that led to the
primary methods. Table~\ref{tab:earlycandidates} uses retrieval-derived
labels, whereas Tables~\ref{tab:qwencheckpoints} and
\ref{tab:secondaryvariants} use whole-question evidence-length labels.
Their macro-F1 columns therefore do not form a single ranking across
targets. The full Phase 2E grid is given in Table~\ref{tab:lrgrid}.

% Generated by scripts/generate_results.py; do not hand-edit.
\begin{table}[htbp]
\centering\small
\caption{All saved question-embedding router candidates.}
\label{tab:earlycandidates}
\begin{tabular}{@{}lrrrrr@{}}
\toprule
Model & LR & Weight decay & Seed & Accuracy & Macro F1\\
\midrule
Logistic & 0.01 & 0 & 42 & 0.2478 & 0.2060\\
Logistic & 0.01 & 0.0001 & 43 & 0.2478 & 0.2149\\
Logistic & 0.001 & 0 & 44 & 0.2511 & 0.2080\\
Logistic & 0.001 & 0.0001 & 45 & 0.2500 & 0.2046\\
MLP & 0.001 & 0 & 10042 & 0.2749 & 0.2187\\
\bottomrule
\end{tabular}
\par\smallskip\begin{minipage}{\textwidth}\footnotesize Legacy retrieval-derived labels. Logistic seed 43 was retained: the MLP's macro-F1 advantage was below the 0.01 adoption threshold. Candidate seeds differ, so these are not seed-controlled hyperparameter ablations.\end{minipage}
\end{table}

% Generated by scripts/generate_results.py; do not hand-edit.
\begin{table}[htbp]
\centering\small
\caption{Three-epoch Qwen checkpoint histories before the Phase 2E grid.}
\label{tab:qwencheckpoints}
\begin{tabular}{@{}lrrrrl@{}}
\toprule
Phase & Epoch & Step & Accuracy & Macro F1 & Selected\\
\midrule
2 & 1 & 71 & 0.4545 & 0.1250 & ---\\
2 & 2 & 142 & 0.4556 & 0.1270 & ---\\
2 & 3 & 213 & 0.4318 & 0.1650 & Yes\\
2B-A & 1 & 71 & 0.2511 & 0.0803 & ---\\
2B-A & 2 & 142 & 0.3009 & 0.1904 & ---\\
2B-A & 3 & 213 & 0.3506 & 0.2092 & Yes\\
2B-B & 1 & 71 & 0.1926 & 0.0646 & ---\\
2B-B & 2 & 142 & 0.3701 & 0.1684 & Yes\\
2B-B & 3 & 213 & 0.2684 & 0.1106 & ---\\
2C & 1 & 71 & 0.4361 & 0.1881 & ---\\
2C & 2 & 142 & 0.2727 & 0.1715 & ---\\
2C & 3 & 213 & 0.3452 & 0.2176 & Yes\\
2D & 1 & 71 & 0.4405 & 0.1807 & ---\\
2D & 2 & 142 & 0.2803 & 0.1678 & ---\\
2D & 3 & 213 & 0.3690 & 0.2299 & Yes\\
\bottomrule
\end{tabular}
\par\smallskip\begin{minipage}{\textwidth}\footnotesize Validation evidence-length classification. Only selected checkpoints have the reported final retrieval evaluation. Phase 2 uses the corrected v2 run; preflights and its interrupted predecessor are not competing full runs.\end{minipage}
\end{table}

% Generated by scripts/generate_results.py; do not hand-edit.
\begin{table}[htbp]
\centering\small
\caption{Secondary similarity and fusion variants on 924 validation questions.}
\label{tab:secondaryvariants}
\begin{tabular}{@{}lrr@{}}
\toprule
Variant & Accuracy & Macro F1\\
\midrule
3A: maximum similarity & 0.0498 & 0.0600\\
3A: top-five mean & 0.0433 & 0.0573\\
3A: penalised top-five & 0.1851 & 0.1434\\
3A: leaf breadth & 0.3755 & 0.1242\\
3A: level-only linear & 0.3236 & 0.1892\\
3B: level-only XGBoost & 0.2933 & 0.1822\\
3C: hidden-state fusion & 0.3203 & 0.2298\\
\bottomrule
\end{tabular}
\par\smallskip\begin{minipage}{\textwidth}\footnotesize These are evidence-length classification diagnostics, not additional retrieval evaluations. Original hidden-state fusion has the same in-sample base-feature limitation as original logits fusion.\end{minipage}
\end{table}


\section{Token-count classifier instruction}
The Phase 2D/2E instruction is reproduced from the repository:
\begin{quote}
You are a router for a retrieval-augmented generation system. Based
only on the question, select the option representing the context
size most suitable for retrieving the evidence required to answer
it. Choose exactly one value from: 1 = 10 tokens, 2 = 20 tokens,
3 = 40 tokens, 4 = 80 tokens, 5 = 160 tokens. Return only the number
\end{quote}
The original question follows a blank line and the prefix
\texttt{Question:}. Although the wording refers to returning a number,
the sequence-classification experiment obtains five head logits,
not a generated response. The prompt is an input specification;
it should not be confused with the model's output mechanism.

\section{Final frozen context-feature instruction}
The repaired Phase 5B-v2 instruction states:
\begin{quote}
You are not being asked to answer the question. Your only task is
to classify how much supporting context would be needed to answer
it. Respond with exactly one lowercase label from: short, medium,
long. Do not output yes, no, the answer to the question, an
explanation, or punctuation. Output one label only.
\end{quote}
The final decision restricts the next-token comparison to
\texttt{short} (8412), \texttt{medium} (25252), and
\texttt{long} (4670). The chosen category is represented by three
one-hot features. It is repeated across local trees for that question.

The free-generation predecessor used a maximum of eight new tokens.
Its invalid responses are retained, and no classifier or retrieval
result is attributed to it. The final constrained method is a
distinct version, not a retroactive replacement of those invalid
records.

\section{Prediction distributions}
% Generated by scripts/generate_results.py; do not hand-edit.
\begin{table}[htbp]
\centering\small
\caption{Predicted question-level sizes on the 924 validation questions.}
\label{tab:predictioncounts}
\begin{tabular}{@{}lrrrrr@{}}
\toprule
Phase & 10 & 20 & 40 & 80 & 160\\
\midrule
1 & 767 & 40 & 116 & 0 & 1\\
2 & 0 & 0 & 0 & 149 & 775\\
2B-A & 0 & 0 & 427 & 189 & 308\\
2B-B & 0 & 0 & 0 & 434 & 490\\
2C & 0 & 20 & 224 & 374 & 306\\
2D & 0 & 16 & 219 & 332 & 357\\
2E & 0 & 15 & 275 & 366 & 268\\
3A & 8 & 31 & 209 & 356 & 320\\
3B & 13 & 32 & 201 & 377 & 301\\
3C & 27 & 88 & 251 & 300 & 258\\
3C-OOF & 9 & 38 & 220 & 385 & 272\\
4 & 4 & 359 & 561 & 0 & 0\\
\bottomrule
\end{tabular}
\end{table}

The prediction counts in Table~\ref{tab:predictioncounts} are useful
for checking both completeness and class collapse. Each row sums
to 924. A row concentrated on one or two sizes can still have a
moderate accuracy under an imbalanced target. It should be inspected
alongside macro F1 and retrieval utility.

\section{Selection and leakage boundaries}
Qwen checkpoints are selected by validation macro F1, then accuracy,
weighted F1, balanced accuracy, lower cross-entropy, and earlier
checkpoint. Phase 2E adds lower learning rate as the final tie-breaker.
Retrieval cannot revise that winner.

Tree-model grids use paper-grouped training folds. Corrected fusion
uses OOF base features, without claiming fully nested meta-level
diagnostics. Phase 4 fixes inherited hyperparameters and saves
validation decisions before opening utility targets. Phases 5A and
5B-v2 similarly separate predictions from gold-overlap evaluation.
Despite these safeguards, validation remains a development split
across the programme as a whole.

\chapter{Repository Provenance and Reproduction}
\label{app:provenance}

\section{How reported numbers map to artifacts}
Paths in this appendix are relative to the project repository,
not to the LaTeX directory. The structured files are the primary
numerical record; the documents in \path{docs/} explain protocol
and execution history.

\begin{description}[style=nextline,leftmargin=0pt]
\item[Data construction and early baseline]
\path{docs/DATA_INTEGRITY.md},
\path{docs/EVALUATION_PIPELINE.md}, and
\path{reports/final_validation_comparison/}.
The selected embedding-router report is
\path{models/granularity_router/frozen_topk5/training_report.json}.
The frozen validation action table is
\path{outputs/oracle_frozen/validation/RouterDataset_20260623_171712.jsonl}.

\item[Question-only Qwen phases]
Each phase directory below contains classification and retrieval
summaries and a final summary:
\path{outputs/qwen_pretrained_zero_shot_router_evidence_length_oracle/},
\path{outputs/qwen_finetuned_router_evidence_length_oracle/},
\path{outputs/qwen_phase2b_alias_unweighted_evidence_length_oracle/},
\path{outputs/qwen_phase2b_alias_classbalanced_evidence_length_oracle/},
\path{outputs/qwen_phase2c_sequence_classifier_evidence_length_oracle/},
and
\path{outputs/qwen_phase2d_sequence_classifier_token_count_prompt_evidence_length_oracle/}.
The Phase 2E study is under
\path{outputs/qwen_phase2e_lr_grid_token_count_prompt_5epochs_evidence_length_oracle/};
the selected retrieval result is under \path{trials/lr5e-6/}.

\item[Similarity features and fusion]
The corresponding roots are
\path{outputs/similarity_tree_phase3a_evidence_length_oracle/},
\path{outputs/similarity_tree_phase3b_xgboost_evidence_length_oracle/},
\path{outputs/qwen_phase3c_fusion_evidence_length_oracle/}, and
\path{outputs/qwen_phase3c_oof_fusion_evidence_length_oracle/}.
Within them, \path{classification_summary.json} records representation
comparisons, while \path{classification/metrics.json} and
\path{retrieval/summary.json} describe the primary result.

\item[Expected regret]
\path{outputs/qwen_phase4_expected_regret_retrieval_utility/comparison/baselines.json}
contains the comparable strategy metrics and paired bootstrap
differences. Its \path{retrieval/results.jsonl} records all five
utilities, the predicted size, and the resulting regret per question.
Its \path{targets/train_utility_targets.jsonl} contains the completed
training utility targets.

\item[Local routing]
\path{outputs/similarity_tree_phase5a_local_gold_overlap_router/retrieval/summary.json}
contains the primary local result, matched fixed-size methods,
the all-descendants probe, and paired intervals.
The construction summaries under \path{features/} document exclusions
and local training units.

\item[Frozen context feature]
The first version is preserved under
\path{outputs/similarity_tree_phase5b_zero_shot_qwen_context_feature/}.
The repaired version is under
\path{outputs/similarity_tree_phase5b_v2_zero_shot_qwen_context_feature/}.
Its \path{comparison/phase5a_and_fixed40.json} records the paired
comparisons, and its \path{development/} and \path{qwen_features/}
directories separate the training-only repair probe from final
feature extraction.
\end{description}

The experiment entry points include \path{granularity_router.py},
\path{router_selected.py}, the \path{qwen_phase*.py} scripts,
\path{similarity_tree_phase3a.py},
\path{similarity_tree_phase3b.py},
\path{similarity_tree_phase5a.py}, and
\path{similarity_tree_phase5b_v2.py}. The phase documents specify
the exact commands and source snapshots used for their runs.
Historical absolute paths in transferred metadata identify the
original GPU host; they are not assumed to exist on the current
writing computer.

\section{Rebuilding the thesis results}
The LaTeX project includes
\path{scripts/generate_results.py}. It reads preserved JSON and JSONL
artifacts and produces the tables in \path{content/generated/} and
vector figures in \path{images/generated/}. It does not call an
embedding service, access the vector database, train a model, or
modify experiment outputs.

The script checks that the main prediction distributions sum to
924, that the Phase 4 records cover 924 questions and 277 papers,
and that per-question fixed-size means reproduce the stored
comparison values. It records SHA-256 hashes of its input files
in \path{content/generated/source_manifest.json}. The diagnostic
cases preserve their question identifiers in that manifest.

A separate read-only check, \path{scripts/audit_results.py}, reconciles
question and paper identities, recomputes classification metrics and
retrieval aggregates from saved rows, compares routed results with
their frozen actions, and reconstructs the reported paper-cluster
bootstrap intervals. It writes its findings to
\path{build/qa/results_audit.json}. This verifies consistency of the
preserved record; it does not independently rerun the original models
or prove that an annotation is semantically complete.

To rebuild those derived assets from the thesis directory:
\begin{verbatim}
python scripts/generate_results.py
python scripts/audit_results.py
\end{verbatim}
The plotting environment needs NumPy and Matplotlib. Generated
tables and figures are retained with the LaTeX source so that
ordinary document compilation does not require the experimental
Python environments.

\section{Compiling and reviewing the document}
The project uses pdfLaTeX and Biber. The included PowerShell
build script runs the required passes and writes the PDF and
intermediate files to \path{build/}. From the thesis directory:
\begin{verbatim}
powershell -ExecutionPolicy Bypass -File scripts/build_thesis.ps1
\end{verbatim}
The existing VS Code LaTeX configuration also provides a build
recipe. Compilation does not require a GPU, a running Qdrant
service, API credentials, or the original model checkpoints.

The original demonstration bibliography and glossary files remain
in the downloaded template as reference assets but are not included
in this thesis. The personal acknowledgements use text supplied by
the author and are included in the front matter.

This document is a complete research draft grounded in the saved
repository. Before submission, the author and supervisors should
review the scientific interpretation, institutional formatting
requirements, and citation details. Any new evaluation should be
added as a separately identifiable result and reflected consistently
in the abstract, summary, results, and conclusion.
